\documentclass[11pt]{article}
\usepackage[margin=1in]{geometry}
\usepackage{amsmath,amssymb,amsthm}
\usepackage{graphicx}
\usepackage[hidelinks]{hyperref}

\newtheorem{theorem}{Theorem}
\newtheorem{definition}{Definition}
\newtheorem{remark}{Remark}

\newcommand{\N}{\mathbb N}
\newcommand{\ellTwo}{\ell_2}
\newcommand{\BM}{d_{\mathrm{BM}}}

\title{Partial Result: A Uniform Tail Property-(H) Obstruction}
\author{Automatic Math Discovery Run \texttt{fa\_banach\_001}}
\date{June 14, 2026}

\begin{document}
\maketitle

\section*{Source Problem}

The source is Valentin Ferenczi, Jordi Lopez-Abad, Behrang Mbombo, and
Sebastien Todorcevic, \emph{Local Banach-space dichotomies and ergodic spaces},
arXiv:2005.06458. Question 6.17 appears on PDF page 60 and asks:
\begin{quote}
1. Does every non-Hilbertian space have a non-Hilbertian subspace with a
Schauder basis?

2. Does every non-Hilbertian space with unconditional FDD have a
non-Hilbertian subspace with an unconditional basis?
\end{quote}
This packet addresses the second question.

\begin{center}
\includegraphics[width=\textwidth]{figures/open_problem_crop.png}
\end{center}

\section*{Definitions}

Let \(X\) be a Banach space with an unconditional finite-dimensional
decomposition (FDD) \((E_n)_{n=1}^\infty\). A normalized block sequence of
\((E_n)\) means a sequence of unit vectors with successive finite supports
with respect to the FDD.

\begin{definition}[uniform tail property-(H) failure along an FDD]
We say that \((E_n)\) has uniform tail property-(H) failure if there is
\(\lambda \ge 1\) such that for every \(N\in\N\) and every \(M\ge1\), there
are integers \(N\le r\le s\) and a finite-dimensional subspace
\[
G\subseteq E_r\oplus E_{r+1}\oplus \cdots \oplus E_s
\]
which has a normalized \(\lambda\)-unconditional basis and satisfies
\[
\BM(G,\ellTwo^{\dim G})>M.
\]
\end{definition}

This condition is a tail version of failing Pisier's property (H) with a fixed
unconditionality constant. It is exactly the finite-dimensional supply used in
the source paper's proof of Theorem 6.15 for \(d_2\)-minimal spaces, but here
it is stated as an explicit hypothesis rather than derived from minimality
among non-Hilbertian spaces.

\section*{Idea of the Proof}

If some normalized block sequence of the FDD already spans a non-Hilbertian
space, the unconditionality of the FDD gives the desired unconditional basis.
Otherwise every normalized block sequence is Hilbertian. A standard gliding
hump argument then upgrades this to a uniform statement: all normalized block
sequences are \(C\)-equivalent to the unit vector basis of \(\ell_2\). The
tail witnesses supply successive finite-dimensional blocks whose internal
bases are uniformly unconditional but whose Hilbertian distortion tends to
infinity. The uniform \(\ell_2\) behavior of normalized block vectors controls
sign changes between the blocks, while the internal unconditionality controls
sign changes inside each block. Concatenating the block bases therefore gives
an unconditional basis. Since the blocks have unbounded Banach-Mazur distance
from Euclidean spaces, their closed span cannot be Hilbertian.

\begin{theorem}
Let \(X\) be a Banach space with an unconditional FDD \((E_n)\). If
\((E_n)\) has uniform tail property-(H) failure, then \(X\) contains a
non-Hilbertian subspace with an unconditional basis.
\end{theorem}

\begin{proof}
Let \(K\) be an unconditionality constant for the FDD. If there is a
normalized block sequence \((u_j)\) of \((E_n)\) whose closed span is
non-Hilbertian, then \((u_j)\) is an unconditional basic sequence, with
constant bounded in terms of \(K\), and we are done.

Assume therefore that every normalized block sequence spans a Hilbertian
subspace. We first record the usual uniformity consequence. There is a
constant \(C\) such that every normalized block sequence of \((E_n)\) is
\(C\)-equivalent to the unit vector basis of \(\ell_2\). Indeed, if no such
constant existed, then for every tail and every integer \(m\) one could choose
a finite normalized block sequence in that tail which is not \(m\)-equivalent
to the Euclidean basis of the same length. Choosing these finite sequences
successively and concatenating them gives a normalized block sequence. By the
current assumption its span is Hilbertian; since the sequence is unconditional,
the standard Hilbert-space fact for normalized unconditional basic sequences
implies that it is equivalent to the \(\ell_2\) basis. This contradicts the
successively chosen finite blocks with constants tending to infinity.

Let \(\lambda\) witness uniform tail property-(H) failure. Inductively choose
successive finite-dimensional block subspaces \(G_k\), each supported after
the previous one, such that \(G_k\) has a normalized
\(\lambda\)-unconditional basis
\[
(x_i)_{i_k\le i<i_{k+1}}
\]
and \(\BM(G_k,\ell_2^{\dim G_k})\ge k\). We claim that the concatenated
sequence \((x_i)_{i=0}^\infty\) is unconditional.

Fix finitely supported scalars \((a_i)\) and signs \((\varepsilon_i)\). For
each \(k\), put
\[
u_k=\sum_{i_k\le i<i_{k+1}} a_i x_i,\qquad
v_k=\sum_{i_k\le i<i_{k+1}} \varepsilon_i a_i x_i.
\]
Let \(b_k=\|u_k\|\) and \(c_k=\|v_k\|\). By the
\(\lambda\)-unconditionality of the basis inside \(G_k\),
\[
c_k\le \lambda b_k.
\]
When \(b_k\ne0\), write \(y_k=u_k/b_k\), and when \(c_k\ne0\), write
\(z_k=v_k/c_k\). Omitting zero terms, \((y_k)\) and \((z_k)\) are normalized
block sequences of the original FDD. Hence they are \(C\)-equivalent to the
\(\ell_2\) basis. Therefore
\[
\left\|\sum_i \varepsilon_i a_i x_i\right\|
=\left\|\sum_k c_k z_k\right\|
\le C\left(\sum_k c_k^2\right)^{1/2}
\le \lambda C\left(\sum_k b_k^2\right)^{1/2}
\le \lambda C^2\left\|\sum_k b_k y_k\right\|
=\lambda C^2\left\|\sum_i a_i x_i\right\|.
\]
Thus \((x_i)\) is unconditional.

Let \(Y=[x_i:i\in\N]\). Since each \(G_k\) is a finite-dimensional subspace of
Y and \(\BM(G_k,\ell_2^{\dim G_k})\to\infty\), the space \(Y\) cannot be
isomorphic to a Hilbert space: a Hilbertian isomorphism with constant \(D\)
would make every finite-dimensional subspace of \(Y\) \(D\)-isomorphic to a
Euclidean space. Hence \(Y\) is a non-Hilbertian subspace of \(X\) with an
unconditional basis.
\end{proof}

\section*{Consequence for Question 6.17(2)}

This is a partial positive result for Question 6.17(2). It shows that any
counterexample with an unconditional FDD must be much closer to Hilbert space
than a naive non-Hilbertian UFDD space: either it has no non-Hilbertian block
subspace, or every fixed unconditionality constant becomes Hilbertian on far
enough tails. In particular, the obstruction to the source question is not
merely the lack of bases inside the finite-dimensional FDD blocks; it must
also prevent uniformly unconditional, increasingly non-Hilbertian finite
blocks from appearing arbitrarily far out.

\section*{Novelty and Search Notes}

The cheap run indexes were searched for arXiv:2005.06458, the title phrase,
``unconditional FDD'', ``non-Hilbertian'', ``property (H)'', ``MNH'', and
related ergodicity keywords. The only relevant existing row was a separate
literature-implied answer about Schlumprecht space ergodicity; no packet for
Question 6.17 or this tail-property route was found. Exact web searches on
June 14, 2026 for the Question 6.17 wording and for property-(H)/UFDD variants
did not surface a separate later answer.

This theorem is not a full solution of Question 6.17(2). It abstracts and
extends the block concatenation mechanism used in the source paper's proof of
Theorem 6.15 from the special \(d_2\)-minimal setting to spaces where the
needed finite-dimensional tail witnesses are assumed directly.

\section*{Human Review Recommendation}

Review the uniformity claim for normalized block sequences and the final
concatenation estimate. The proof is short and should be checked mainly for
whether the stated tail witness hypothesis is the right level of generality
for a useful partial result.

\begin{thebibliography}{9}
\bibitem{FLMT}
V. Ferenczi, J. Lopez-Abad, B. Mbombo, and S. Todorcevic,
\emph{Local Banach-space dichotomies and ergodic spaces},
arXiv:2005.06458.
\end{thebibliography}

\end{document}
